\chapter{Анализ предметной области и существующих решений}

\section{Биологические последовательности как объект сонификации}

Биологическая последовательность представляет собой строку символов, описывающую структуру ДНК, РНК или белка. В ДНК используются четыре основания: аденин (A), цитозин (C), гуанин (G) и тимин (T). В РНК тимин заменён на урацил (U). Белковые последовательности описываются алфавитом из 20 аминокислот.

Несмотря на строковую форму, биологические последовательности являются сложными объектами. Они содержат повторы, участки с различным GC-составом, потенциальные кодирующие области, фрагменты с разной энтропией и статистикой k-mer паттернов. Геном человека содержит около 3 миллиардов пар оснований, что делает невозможным непосредственное восприятие последовательности как текста.

Традиционные методы анализа включают визуализацию (графики GC-content, тепловые карты), статистический анализ и вычислительные методы. Сонификация предлагает альтернативный способ представления, использующий слуховое восприятие временных структур.

\subsection{Почему музыка подходит для представления биологических данных}

Музыка обладает несколькими свойствами, делающими её подходящей для сонификации:

\textbf{Временная структура.} Музыка разворачивается во времени, как и биологическая последовательность. Это создаёт естественную основу для отображения линейных данных.

\textbf{Многомерность.} Музыкальная фраза одновременно характеризуется высотой, длительностью, темпом, плотностью, регистром и гармонией. Биологический фрагмент также описывается набором признаков, что позволяет отобразить несколько измерений одновременно.

\textbf{Целостное восприятие.} Человек может услышать, что один фрагмент звучит напряжённее, плотнее или стабильнее другого, даже не зная точных численных значений.

\subsection{Ограничения музыкальной интерпретации}

Важно понимать границы сонификации. Если из ДНК получилась мелодия, это не означает, что мелодия «содержалась» в ДНК. Корректнее говорить, что из биологической последовательности извлекаются признаки, которые управляют генератором музыки. Музыка является результатом модели, обученной на музыкальном корпусе, а не прямым звуковым отпечатком природы.

\section{Информационно-теоретическое обоснование}

Предварительный анализ, выполненный перед разработкой системы, показал, что биологические и музыкальные данные удобно рассматривать как формальные языки. В обоих случаях имеется конечный алфавит, последовательная структура, статистические закономерности и несколько уровней абстракции. Для биологических данных такими уровнями являются нуклеотиды, кодоны, аминокислоты и белковые домены. Для музыки — отдельные MIDI-события, ноты, интервалы, аккорды, такты и фразы.

С точки зрения теории информации~\cite{shannon1948} важно не только количество возможных символов, но и фактическое распределение этих символов. Максимальная энтропия нуклеотидного алфавита равна
\begin{equation}
H_{DNA}^{max} = \log_2 4 = 2
\end{equation}
битам на символ. Для кодонов алфавит состоит из 64 триплетов, поэтому
\begin{equation}
H_{codon}^{max} = \log_2 64 = 6
\end{equation}
бит на кодон. Для аминокислотного алфавита из 20 символов максимальная энтропия составляет
\begin{equation}
H_{aa}^{max} = \log_2 20 \approx 4.32
\end{equation}
бита на символ.

Музыкальные последовательности также можно кодировать на разных уровнях. Событийное MIDI-представление содержит высоты, длительности, velocity, дорожки и временные смещения, но само по себе плохо выражает гармоническую семантику. Более высокоуровневые представления — например, аккорды, относительные ступени, позиции внутри такта и токены музыкальной структуры — имеют более крупный алфавит, но меньшую фактическую случайность, потому что музыка подчиняется метрическим, ладовым и гармоническим ограничениям~\cite{muller2015}.

Из этого следует важный методологический вывод: прямое отображение нуклеотидов в ноты недостаточно, потому что оно сопоставляет разные уровни абстракции. Один нуклеотид несёт слишком мало контекстной информации для выбора музыкального события, а отдельная нота без гармонии и ритма не описывает музыкальную структуру. Поэтому более корректно сначала извлекать признаки биологического фрагмента, а затем использовать их как управляющий сигнал для генерации музыкальных структур.

Эту идею можно выразить через взаимную информацию:
\begin{equation}
I(B;M) = H(M) - H(M \mid B),
\end{equation}
где \(B\) — биологический фрагмент, \(M\) — музыкальный результат. Цель системы состоит не в сохранении всей информации исходной последовательности, а в сохранении структурно значимых признаков, которые могут управлять темпом, плотностью, регистром, модальностью и гармонической сложностью. Такой подход делает сонификацию не буквальным переводом символов, а управляемым сжатием и музыкальной интерпретацией биологических данных.

\section{Обзор существующих подходов к биосонификации}

\subsection{Прямое отображение символов}

Самый простой подход — назначить каждому символу биологического алфавита музыкальную ноту. Например, A→C, C→E, G→G, T→B. Такой метод легко реализовать и объяснить.

\textbf{Преимущества:}
\begin{itemize}
\item Простота реализации
\item Прозрачность отображения
\item Не требует обучения
\end{itemize}

\textbf{Недостатки:}
\begin{itemize}
\item Не использует биологическую информацию (состав, энтропия, структура)
\item Музыкально бедный результат
\item Отсутствие гармонической структуры
\item Монотонность при длинных последовательностях
\end{itemize}

Примером является работа Ohno и Ohno~\cite{ohno1986}, где гены преобразовывались в мелодии путём прямого отображения кодонов.

\subsection{Алгоритмическая композиция}

Алгоритмический подход использует правила композиции, параметры которых определяются биологическими признаками. Например, GC-content может влиять на темп, энтропия — на плотность нот, а длина повторов — на ритмическую структуру.

\textbf{Преимущества:}
\begin{itemize}
\item Учитывает биологические признаки
\item Контролируемое отображение
\item Музыкально более богатый результат
\end{itemize}

\textbf{Недостатки:}
\begin{itemize}
\item Требует ручной настройки правил
\item Ограниченная выразительность
\item Сложность масштабирования
\end{itemize}

Примером является система DNA Music~\cite{temple2017}, использующая параметрическую сонификацию биологических признаков.

\subsection{Генеративные модели без структурного разделения}

Современные подходы используют нейронные сети для генерации музыки. Однако прямая генерация всех музыкальных событий одновременно создаёт проблемы: модель должна одновременно выбирать аккорды, мелодические ноты и организовывать форму.

\textbf{Преимущества:}
\begin{itemize}
\item Обучение на реальных музыкальных данных
\item Потенциально высокая выразительность
\end{itemize}

\textbf{Недостатки:}
\begin{itemize}
\item Хаотичный результат на небольших корпусах
\item Отсутствие музыкальной структуры
\item Сложность управления от биологических признаков
\end{itemize}

\section{Обзор методов генерации символической музыки}

\subsection{Рекуррентные нейронные сети}

Ранние работы использовали LSTM для генерации музыки~\cite{eck2002}. RNN хорошо моделируют последовательности, но имеют ограничения по длине контекста.

\subsection{Вариационные автоэнкодеры}

MusicVAE~\cite{roberts2018} использует VAE для генерации музыкальных фраз. Модель обучается кодировать музыку в латентное пространство, что позволяет интерполировать между композициями.

\subsection{Генеративно-состязательные сети}

MuseGAN~\cite{dong2018} использует GAN для генерации многодорожечной музыки. Модель генерирует несколько инструментальных партий одновременно.

Генеративно-состязательная схема была рассмотрена как один из возможных вариантов для данной работы, поскольку она естественно соответствует идее ``генератор + оценщик'': генератор создаёт музыкальный фрагмент, а дискриминатор или другая модель-оценщик определяет, насколько результат похож на целевое распределение. Однако для дипломной постановки такой подход имеет существенные риски. GAN сложнее стабилизировать на небольших корпусах, труднее интерпретировать численно, а качество обучения сильно зависит от баланса генератора и дискриминатора. Для задачи, где важны воспроизводимость, ограниченный GPU и понятная связь биологических признаков с музыкальными параметрами, более оправданным оказался условный Transformer с внешней экспериментальной оценкой.

\subsection{Трансформеры}

Music Transformer~\cite{huang2018} применяет архитектуру трансформера для генерации длинных музыкальных последовательностей. Механизм внимания позволяет моделировать долгосрочные зависимости.

\section{Сравнение существующих решений}

В таблице~\ref{tab:solutions} представлено сравнение рассмотренных подходов.

\begin{longtable}{|p{0.2\textwidth}|p{0.3\textwidth}|p{0.4\textwidth}|}
\caption{Сравнение подходов к биосонификации}
\label{tab:solutions}\\
\hline
\textbf{Подход} & \textbf{Преимущества} & \textbf{Недостатки} \\
\hline
\endfirsthead
\hline
\textbf{Подход} & \textbf{Преимущества} & \textbf{Недостатки} \\
\hline
\endhead
\hline
\endfoot
Прямое отображение & Простота, прозрачность & Музыкально бедный результат, нет структуры \\
\hline
Алгоритмическая композиция & Учёт признаков, контроль & Ручная настройка, ограниченная выразительность \\
\hline
RNN & Обучение на данных & Ограничение контекста \\
\hline
VAE & Латентное пространство & Сложность управления \\
\hline
GAN & Многодорожечность & Нестабильность обучения \\
\hline
Transformer & Долгосрочные зависимости & Требует больших данных \\
\hline
\end{longtable}

\section{Недостатки существующих решений}

Анализ существующих подходов выявил следующие недостатки:

\begin{enumerate}
\item \textbf{Отсутствие структурного разделения.} Большинство генеративных моделей создают музыку как единый поток событий, что приводит к хаосу на небольших корпусах.

\item \textbf{Слабая связь с биологическими признаками.} Прямое отображение символов не использует биологическую информацию, а сложные модели плохо управляются внешними признаками.

\item \textbf{Отсутствие воспроизводимых экспериментов.} Многие работы не предоставляют код, данные и метрики для проверки результатов.

\item \textbf{Музыкальная несвязность.} Генерация без учёта гармонической структуры даёт музыкально бедный результат.
\end{enumerate}

Отдельно следует отметить ограничение генеративно-состязательных подходов для рассматриваемой задачи. Хотя модель-оценщик могла бы использоваться как часть обучающего цикла, такой контур усложняет доказательство результата: улучшение функции дискриминатора не всегда означает улучшение музыкальной связности или управляемости от биологических признаков. Поэтому в данной работе оценщик не включается в процесс обучения как в GAN, а используется после генерации в виде набора метрик качества.

\section{Постановка задачи}

На основе анализа предметной области и существующих решений формулируется следующая задача:

\textbf{Разработать систему генерации символической музыки, которая:}
\begin{itemize}
\item Принимает на вход биологические последовательности (DNA, RNA, protein)
\item Извлекает биологические признаки (состав, энтропия, белковые свойства)
\item Использует иерархическую архитектуру: сначала генерирует гармонию, затем мелодию
\item Обучается на реальном музыкальном корпусе
\item Создаёт музыкально связные двухдорожечные MIDI-файлы
\item Обеспечивает воспроизводимость результатов
\end{itemize}

\textbf{Требования к системе:}
\begin{itemize}
\item Поддержка FASTA-формата
\item Обучение на доступном GPU (RTX 2060 6GB)
\item Экспериментальная оценка качества
\item Сравнение с baseline
\item Документация и примеры использования
\end{itemize}

\section*{Выводы по главе 1}

В первой главе проанализирована предметная область биосонификации и генерации символической музыки. Рассмотрены существующие подходы: прямое отображение символов, алгоритмическая композиция и генеративные модели. Выявлены недостатки: отсутствие структурного разделения, слабая связь с биологическими признаками, музыкальная несвязность.

Сформулирована задача разработки системы с иерархической архитектурой Bio→Harmony→Melody, которая генерирует музыкально связные композиции под управлением биологических признаков.
